% ─────────────────────────────────────────────────────────────────────────────
% Overleaf free-tier compatibility test
%
% Tests:
%   1. backend=bibtex (lighter than biber, common on Overleaf)
%   2. backend=biber  (standard)
%   3. style=numeric  (no compression — plain style)
%   4. style=numeric-comp (with compression)
%   5. pdflatex engine (Overleaf default)
%
% Switch the biblatex line below to test each combination.
% ─────────────────────────────────────────────────────────────────────────────
\documentclass[11pt, a4paper]{article}

\usepackage{lmodern}
\usepackage[margin=2.5cm]{geometry}

%%% COMBINATION TO TEST (uncomment one):
%% A) numeric-comp + biber (full compression, standard)
\usepackage[style=numeric-comp, backend=biber, sorting=none]{biblatex}
%% B) numeric-comp + bibtex (compression, lighter backend)
%\usepackage[style=numeric-comp, backend=bibtex, sorting=none]{biblatex}
%% C) numeric + biber (no compression, standard backend)
%\usepackage[style=numeric, backend=biber]{biblatex}
%% D) numeric + bibtex (no compression, lighter backend)
%\usepackage[style=numeric, backend=bibtex]{biblatex}

\begin{filecontents}{overleaf-refs.bib}
@article{ref1, author={Einstein, Albert}, title={Title One}, journal={J. Phys.}, year={2001}}
@article{ref2, author={Bohr, Niels}, title={Title Two}, journal={J. Phys.}, year={2002}}
@article{ref3, author={Dirac, Paul}, title={Title Three}, journal={J. Phys.}, year={2003}}
@article{ref4, author={Feynman, Richard}, title={Title Four}, journal={J. Phys.}, year={2004}}
@article{ref5, author={Heisenberg, Werner}, title={Title Five}, journal={J. Phys.}, year={2005}}
@book{ref6, author={Wald, Robert M.}, title={General Relativity}, publisher={Chicago}, year={1984}}
@book{ref7, author={Misner, C. W. and Thorne, K. S. and Wheeler, J. A.}, title={Gravitation}, publisher={Freeman}, year={1973}}
@article{ref8, author={Hawking, Stephen W.}, title={Title Eight}, journal={Nature}, year={2008}}
\end{filecontents}
\addbibresource{overleaf-refs.bib}

\usepackage{citemsp}
\usepackage[colorlinks=true, citecolor=blue]{hyperref}

\begin{document}

\section*{Overleaf Compatibility Test}

\subsection*{Plain citations (range compression if numeric-comp)}

\noindent
\texttt{\textbackslash cite}: \cite{ref1,ref2,ref3,ref4,ref5} \\
\texttt{\textbackslash citemsp}: \citemsp{ref1, ref2, ref3, ref4, ref5} \\
Non-consecutive: \citemsp{ref1, ref2, ref3, ref5, ref6, ref7}

\subsection*{Locator entries}

\noindent
Section:     \citemsp{ref6/6.1} \\
Sec + para:  \citemsp{ref6/6.1/2} \\
Equation:    \citemsp{ref6/3.2/e4.3} \\
Figure:      \citemsp{ref6/6.1/f3} \\
Footnote:    \citemsp{ref8/1/n7} \\
Page:        \citemsp{ref6/p42} \\
Table:       \citemsp{ref7/8.4/t2} \\
Definition:  \citemsp{ref6/3.2/d5}

\subsection*{Mixed: ranges + locators}

\noindent
Range then locator:  \citemsp{ref1, ref2, ref3, ref4/3.2} \\
Locator then range:  \citemsp{ref1/3.2, ref2, ref3, ref4, ref5} \\
Interleaved:         \citemsp{ref1, ref2, ref3, ref4/e2.1, ref5, ref6, ref7}

\subsection*{Italic context}

\textit{The derivation~\citemsp{ref6/3.2/e4.3} proceeds from the Bianchi identity.}

\subsection*{Realistic paragraph}

The Schwarzschild metric~\citemsp{ref6/6.1/2} describes the vacuum exterior of
a spherically symmetric mass. The geodesic equation can be found in standard
references~\citemsp{ref6/3.2/e4.3, ref7/p520}. For the original field
equations see~\citemsp{ref1, ref2, ref3, ref4, ref5}.

\printbibliography

\end{document}
